\chapter{Orquestador y Coordinación de Flota Multi-Robot}
\label{chap:orquestador}

\section{Arquitectura del orquestador central en C++20}
\label{sec:orquestador_c20}

El orquestador de flota actúa como el puesto de mando centralizado que supervisa el estado global de la brigada de rescate. Diseñado en C++20, centraliza la recepción de telemetría de los tres cuadrúpedos y gestiona la asignación óptima de misiones.

\section{Mecanismo de despacho de órdenes (ROS 2 Actions)}
\label{sec:despacho_ordenes}

La interacción entre el orquestador global y los árboles de comportamiento individuales se implementa mediante interfaces de acción estándar de ROS 2 (\texttt{ExecuteTask.action}):
\begin{itemize}
    \item \textbf{Meta (\textit{Goal}):} Identificador de la tarea, tipo de misión (\texttt{task\_type}) y datos geométricos (objetivo o lista de waypoints).
    \item \textbf{Retroalimentación (\textit{Feedback}):} Porcentaje de avance, posición instantánea del robot y estado de salud operativa.
    \item \textbf{Resultado (\textit{Result}):} Confirmación de éxito (\texttt{SUCCESS}), aborto por contingencia o reporte de causas de fallo.
\end{itemize}

\section{Sincronización y prevención de conflictos entre agentes}
\label{sec:sincronizacion_agentes}

\textit{[Sección esquemática: Reglas de prioridad para evitar bloqueos mutuos en pasillos y coordinación para misiones de relevo y marcaje cooperativo].}
